\section{An exact elliptic gate for a nonidentity cubic unit}
\label{eqg-section}

The unit \(u=\zeta\) at \(g=3\) remains after the identity-unit
exclusion. Its first quotient gives a concrete genus-two curve
and an elliptic lifting condition. The following ordinary
proofs use no software rank output.

\begin{theorem}[Two explicit elliptic quotients]\label{eqg-quotients}
Let \(H\) be the smooth projective genus-two curve
\begin{equation}\label{eqg-sextic}
 y^2=P(s),\qquad
 P(s)=s^6+3s^5-3s^4-11s^3-3s^2+3s+1.
\end{equation}
This is the first HG quotient for \(g=3,u=\zeta\).
It has degree-two morphisms to
\begin{equation}\label{eqg-elliptic-curves}
 E:\ Y^2=X^3-9X-9,\qquad
 E':\ (Y')^2=(X')^3-189X'+999,
\end{equation}
given on finite charts by
\begin{align}
 X&=-\frac{2s^2+5s+2}{(s+1)^2},
 &Y&=\frac{y}{(s+1)^3},\label{eqg-first-map}\\
 X'&=\frac{3(2s^2-7s+2)}{(s-1)^2},
 &Y'&=-\frac{9y}{(s-1)^3}.\label{eqg-second-map}
\end{align}
Moreover \(\operatorname{Jac}(H)\) is isogenous over
\(\mathbb Q\) to \(E\times E'\).
\end{theorem}
\begin{proof}
The identity-unit cubic coordinates are
\(A_0=s^3-3s-1\), \(B_0=3s(s+1)\).
Multiplication by \(\zeta\) gives \(A=-B_0\), \(B=A_0+B_0\);
thus \(B(A+B)=A_0(A_0+B_0)=P\).
On the chart \(s\ne1\), put
\[
 t=\frac{s+1}{s-1},\qquad v=y(t-1)^3.
\]
Direct expansion yields
\begin{equation}\label{eqg-even-sextic}
 v^2=-9t^6+99t^4-27t^2+1.
\end{equation}
The inverse is \(s=(t+1)/(t-1)\), \(y=v/(t-1)^3\);
the birational change extends to the smooth projective models.
In this chart the elliptic coordinates are
\begin{align*}
 X&=(1/t^2-9)/4,&Y&=v/(8t^3),\\
 X'&=(33-9t^2)/4,&Y'&=-9v/8.
\end{align*}
Substitution proves the equations and
\eqref{eqg-first-map}--\eqref{eqg-second-map}.
The two discriminants are \(11664\) and \(944784\), respectively.
The second map is the quotient by \(\sigma:(t,v)\mapsto(-t,v)\);
the first is the quotient by \(\sigma\iota:(t,v)\mapsto(-t,-v)\),
where \(\iota\) is the hyperelliptic involution. The indicated
invariants together with \(t^2\) generate the quotient fields.
The involutions are nontrivial, so both maps have degree two,
and extend over the smooth projective models.

Let \(\Phi\) be the pair of quotient norms and \(\Psi\) the sum
of pullbacks. The hyperelliptic quotient is \(\mathbb P^1\);
its divisor norm/pullback identity gives \(\iota=-1\) on
\(\operatorname{Jac}(H)\). The other two double covers therefore give
\begin{equation}\label{eqg-times-two}
 \Psi\Phi=(1+\sigma)+(1+\sigma\iota)=[2].
\end{equation}
These divisor identities include the ramified fibres.
Multiplication by two is surjective, hence so is \(\Psi\).
Both sides have dimension two, so \(\Psi\) is a
\(\mathbb Q\)-isogeny. This is the explicit arithmetic
norm/pullback argument of Theorem~\ref{hg-isogeny}.
\end{proof}

\begin{theorem}[The complete rational image of the first quotient]
\label{eqg-image}
Under \eqref{eqg-first-map}, the image of \(H(\mathbb Q)\) is exactly
\begin{equation}\label{eqg-square-gate}
 \{O\}\ \cup\
 \{(X,Y)\in E(\mathbb Q):4X+9\text{ is a square in }\mathbb Q\}.
\end{equation}
\end{theorem}
\begin{proof}
Clearing the abscissa formula gives
\begin{equation}\label{eqg-inverse-quadratic}
 (X+2)s^2+(2X+5)s+(X+2)=0,
\end{equation}
with discriminant \(4X+9\). At \(s=-1\), \(P(-1)=1\),
and the two rational points map to \(O\). The two rational
points at infinity, and those at \(s=0\), map to
\((-2,1)\), \((-2,-1)\). Thus every chart exception is
also in \eqref{eqg-square-gate}.

Conversely, let a finite elliptic point have \(4X+9\) square.
If \(X=-2\), then \(Y=\pm1\), and \(s=0,y=Y\) is a preimage.
Otherwise \eqref{eqg-inverse-quadratic} has a rational root \(s\).
This root cannot be \(-1\), since the left side at \(-1\)
is \(-1\). Put \(y=Y(s+1)^3\). The polynomial identity
behind \eqref{eqg-first-map} gives \(y^2=P(s)\).
The repeated-discriminant case cannot occur on \(E(\mathbb Q)\):
at \(X=-9/4\) its cubic has value \(-9/64\).
Finally \(O\) has the two rational preimages already listed.
\end{proof}

This is one quotient condition for the actual seed problem.
A point of \(D_{3,\zeta}\) must also have the second HG square
root at the same source \(s\), with \(0<A(s)/B(s)<1/3\).
Arbitrary elliptic points, or quotient points at different
source parameters, do not construct an actual seed.

\begin{theorem}[Two independent rank contributions]\label{eqg-rank-lower}
Both elliptic curves in \eqref{eqg-elliptic-curves} have points
of infinite order. Thus \(\operatorname{Jac}(H)(\mathbb Q)\)
has Mordell--Weil rank at least two.
\end{theorem}
\begin{proof}
Take \(P=(-2,1)\in E(\mathbb Q)\) and \(P'=(6,9)\in E'(\mathbb Q)\).
Direct doubling gives
\[
 2P=(25/4,-107/8),\qquad 2P'=(33/4,9/8).
\]
For a short integral equation \(y^2=x^3+ax+b\) with \(a,b\) odd,
\[
 x(2Q)=\frac{x^4-2ax^2-8bx+a^2}{4(x^3+ax+b)}.
\]
If \(v_2(x)=s<0\), the unique smallest numerator and denominator
valuations are \(4s\) and \(2+3s\). Hence
\(v_2(x(2Q))=s-2\). Both displayed doubled points start at
valuation \(-2\), and their repeated doublings have distinct
abscissas with valuations \(-2,-4,-6,\ldots\).
Both original points therefore have infinite order.
The \(\mathbb Q\)-isogeny of Theorem~\ref{eqg-quotients} and
the Mordell--Weil theorem give the rank lower bound.
No assertion that these points generate the full groups is used.
\end{proof}

Consequently a rank-less-than-genus argument for this genus-two
quotient cannot be justified. Other descents and simultaneous
square conditions remain available. The gate itself is strict:
\(P\) lifts but \(2P\) does not, since \(4(25/4)+9=34\) is
not a rational square. This makes no claim about all higher multiples.

The separate bounded PARI/GP search, checked by exact rational
arithmetic, found \(24\) affine points across the nine cubic
quotient models with numerator and denominator bounded by \(1000\).
All had \(s=0\) or \(-1\), and none supplied the positive
common-source locus. A distinct exhaustive rational-abscissa check
at bound \(80\) checked \(70767\) curve-abscissa cases.
The same replay checked all \(648\) primitive coordinate classes
modulo \(27\) for Theorem~\ref{cl-no-local-point}.
These finite checks do not determine all rational points.
Diagnostic software rank outputs are not used in the ordinary
proofs above, and no genus-two or elliptic-curve theorem here is
claimed as a complete Lean formalization.
